\documentclass[12pt,a4paper]{article}

% -------------------------------------------------
% Sprache / Kodierung
% -------------------------------------------------
\usepackage[T1]{fontenc}
\usepackage[utf8]{inputenc}
\usepackage[ngerman]{babel}
\usepackage{lmodern}
\usepackage{microtype}

% -------------------------------------------------
% Mathematik
% -------------------------------------------------
\usepackage{amsmath,amssymb,amsthm,mathtools}

% -------------------------------------------------
% Layout / Grafik / Farben
% -------------------------------------------------
\usepackage[a4paper,margin=2.3cm]{geometry}
\usepackage{booktabs}
\usepackage{xcolor}
\usepackage[hidelinks,unicode]{hyperref}
\pdfstringdefDisableCommands{%
  \def\ü{ue}\def\ö{oe}\def\ä{ae}\def\Ü{Ue}\def\Ö{Oe}\def\Ä{Ae}\def\ss{ss}%
}
\usepackage[most]{tcolorbox}
\usepackage{enumitem}
\usepackage{array}

% -------------------------------------------------
% Farben (konsistent mit Vorgängerdokumenten)
% -------------------------------------------------
\definecolor{lückenrot}{RGB}{180,40,40}
\definecolor{bewiesengrün}{RGB}{30,100,50}
\definecolor{warnorange}{RGB}{200,100,10}
\definecolor{hintergrundblau}{RGB}{235,242,250}
\definecolor{hintergrundgelb}{RGB}{255,252,225}
\definecolor{adischviolett}{RGB}{90,20,140}
\definecolor{fehlerrot}{RGB}{200,0,0}
\definecolor{konjviolett}{RGB}{80,0,130}
\definecolor{e8grau}{RGB}{55,55,90}
\definecolor{korrekturorange}{RGB}{200,100,10}

% -------------------------------------------------
% tcolorbox-Stile
% -------------------------------------------------
\tcbset{
  lücke/.style={colback=red!8, colframe=lückenrot, fonttitle=\bfseries,
    title={Offene Lücke}},
  ergebnis/.style={colback=green!7, colframe=bewiesengrün,
    fonttitle=\bfseries, title={Bewiesenes Ergebnis}},
  warnung/.style={colback=orange!10, colframe=warnorange,
    fonttitle=\bfseries, title={Achtung / Heuristik}},
  adisch/.style={colback=violet!6, colframe=adischviolett,
    fonttitle=\bfseries, title={2-adische Perspektive}},
  kernidee/.style={colback=hintergrundblau, colframe=blue!60!black,
    fonttitle=\bfseries, title={Kernidee}},
  konj/.style={colback=violet!5, colframe=konjviolett, fonttitle=\bfseries,
    title={Konjektur (unbewiesen)}},
  lte/.style={colback=yellow!8, colframe=e8grau, fonttitle=\bfseries,
    title={LTE-Analyse}},
  hauptsatz/.style={colback=green!10, colframe=bewiesengrün!80!black,
    fonttitle=\bfseries\large, title={Stärkster beweisbarer Satz}},
  korrektur/.style={colback=yellow!5, colframe=orange!70!black,
    fonttitle=\bfseries},
}

% -------------------------------------------------
% Theorem-Umgebungen
% -------------------------------------------------
\newtheorem{definition}{Definition}[section]
\newtheorem{proposition}[definition]{Proposition}
\newtheorem{theorem}[definition]{Theorem}
\newtheorem{lemma}[definition]{Lemma}
\newtheorem{korollar}[definition]{Korollar}
\newtheorem{bemerkung}[definition]{Bemerkung}
\newtheorem{hypothese}[definition]{Hypothese}
\newtheorem{satz}[definition]{Satz}

\newenvironment{beweis}{\begin{proof}[Beweis]}{\end{proof}}

% -------------------------------------------------
% Makros
% -------------------------------------------------
\newcommand{\vii}{\nu_2}
\newcommand{\U}{\mathcal{U}}
\newcommand{\N}{\mathbb{N}}
\newcommand{\Z}{\mathbb{Z}}
\newcommand{\R}{\mathbb{R}}
\newcommand{\Ztwo}{\mathbb{Z}_2}
\newcommand{\abs}[1]{\lvert#1\rvert}
\newcommand{\atwo}[1]{\lvert#1\rvert_2}

% -------------------------------------------------
% Dokument
% -------------------------------------------------
\title{%
  Mehrblock-Drift und das Scheitern der lokalen Kontraktion:\\[0.3em]
  Vom Einzelblock-Argument zur globalen Dynamik\\[0.5em]
  \large Ergänzung zur Collatz-Synthese vom 15.\ Juni 2026%
}
\author{Thomas Hoffbauer}
\date{15.\ Juni 2026}

\begin{document}
\maketitle

\begin{abstract}
Dieses Dokument ergänzt \textit{collatz\_schlussartikel.tex} um eine zentrale
neue Erkenntnis: Die lineare DC-Schranke $\nu \gtrsim 0{,}585\,k$, die einen
einzelnen C$^k$A-Block zur Kontraktion zwingen würde, ist der
\textbf{falsche Mechanismus} --- und das lässt sich nun \emph{beweisen}.

Für extremale C-Ketten mit $n_0+1 = 2^{k+1}\cdot 3^r$ (Startpunkte der
Familie $m = 3^r$) liefert das LTE-Lemma $\nu = 3 + \vii(k+r+1) = O(\log k)$.
Das reicht für Einzelblock-Kontraktion \emph{prinzipiell nicht} aus.

Das richtige Ziel ist ein \emph{Mehrblock-Drift-Argument}: Nicht jeder
einzelne Block $\mathrm{C}^{l_j}\mathrm{A}$ muss kontrahieren, sondern das
Produkt der Nettofaktoren muss asymptotisch gegen~$0$ gehen. Dies ist
strukturell der Ansatz von Tao (2022) --- globale Drift statt lokaler
Kontraktion.
\end{abstract}

\tableofcontents

\bigskip
\begin{tcolorbox}[warnung, title={Einordnung}]
Die Collatz-Vermutung ist unbewiesen. Dieses Dokument zeigt, warum ein
konkreter Beweisansatz (Einzelblock-Kontraktion) fundamental scheitert,
und präzisiert das richtige Beweisziel. Das ist eine \textbf{Verbesserung
des Kenntnisstands}, keine Niederlage.
\end{tcolorbox}

% ====================================================
\section{Das Scheitern der lokalen Kontraktion}
\label{sec:scheitern}
% ====================================================

\subsection{Die Extremfamilie $m = 3^r$}

Wir betrachten Startwerte der Form
\[
  n_0 + 1 = 2^{k+1} \cdot 3^r, \quad r \geq 1,
\]
also $n_0 = 2^{k+1}\cdot 3^r - 1$. Diese Werte starten eine C-Kette der
maximalen Länge~$k$ (denn $\vii(n_0+1) = k+1$, und die Kettenlänge ist
$\leq \vii(n_0+1) - 1 = k$).

Nach der C-Ketten-Formel gilt:
\[
  n_k + 1 = \frac{3^k(n_0+1)}{2^k} = 3^k \cdot 2 \cdot 3^r = 2\cdot 3^{k+r}.
\]
Also:
\[
  n_k = 2\cdot 3^{k+r} - 1.
\]

\subsection{LTE-Berechnung des A-Schritts}

\begin{theorem}[Extremale Tiefe des EA-Schritts]
\label{thm:lte-extrem}
Für $n_0+1 = 2^{k+1}\cdot 3^r$ gilt am Ende der C-Kette:
\[
  \vii(3n_k + 1) = 1 + \vii(3^{k+r+1} - 1).
\]
Konkret:
\[
  \vii(3n_k+1) =
  \begin{cases}
    2 & \text{falls $k+r+1$ ungerade,}\\
    3 + \vii(k+r+1) & \text{falls $k+r+1$ gerade.}
  \end{cases}
\]
Insbesondere: $\vii(3n_k+1) = O(\log k)$ für $k \to \infty$.
\end{theorem}

\begin{beweis}
Es ist $3n_k + 1 = 3(2\cdot 3^{k+r}-1)+1 = 6\cdot 3^{k+r} - 2 = 2(3^{k+r+1}-1)$.
Daher $\vii(3n_k+1) = 1 + \vii(3^{k+r+1}-1)$.
Das LTE-Lemma für $p=2$ und die Basis~$3$ (mit $2 \mid 3-1$) liefert:
\[
  \vii(3^N - 1) =
  \begin{cases}
    1 & \text{falls $N$ ungerade,}\\
    \vii(N) + 2 & \text{falls $N$ gerade.}
  \end{cases}
\]
Mit $N = k+r+1$ folgt die Behauptung. Da $\vii(k+r+1) = O(\log k)$, gilt
auch $\vii(3n_k+1) = O(\log k)$.
\end{beweis}

\subsection{Divergenz des Nettofaktors}

\begin{theorem}[Lokale Kontraktion scheitert fundamental]
\label{thm:kein-einzelblock}
Für $n_0 + 1 = 2^{k+1}\cdot 3^r$ (extremale Familie) gilt:
\[
  \frac{n_{k+1}}{n_0}
  \approx \frac{(3/2)^{k+1}}{2^{\vii(3n_k+1)}}
  \geq \frac{(3/2)^{k+1}}{2^{3+\vii(k+r+1)}}
  \;\xrightarrow{k\to\infty}\; +\infty.
\]
Der Nettofaktor des C$^k$A-Blocks divergiert. Einzelne Blöcke kontrahieren
im Extremfall \textbf{nicht} für große~$k$.
\end{theorem}

\begin{beweis}
Nach dem A-Schritt gilt:
$n_{k+1} = (3n_k+1)/2^{\nu}$ mit $\nu = \vii(3n_k+1)$.
Damit:
\[
  \frac{n_{k+1}}{n_0}
  = \frac{(3^{k+1}m-1)/2^{\nu-1}}{2^{k+1}\cdot m - 1}
  \approx \frac{3^{k+1}m}{2^{\nu-1} \cdot 2^{k+1} m}
  = \frac{3^{k+1}}{2^{\nu+k}}.
\]
Für $\nu = 3 + \vii(k+r+1)$ und $m = 3^r$:
\[
  \frac{n_{k+1}}{n_0}
  \approx \frac{3}{2^3} \cdot \frac{(3/2)^k}{2^{\vii(k+r+1)}}
  = \frac{3}{8} \cdot \frac{(3/2)^k}{2^{\vii(k+r+1)}}.
\]
Da $(3/2)^k$ exponentiell wächst und $2^{\vii(k+r+1)} \leq k+r+1$ nur
linear, divergiert der Quotient.
\end{beweis}

\begin{tcolorbox}[ergebnis, title={Nicht-Ergebnis als Ergebnis}]
\textbf{Theorem~\ref{thm:kein-einzelblock} ist ein bewiesenes Nicht-Ergebnis.}
Es belegt, dass kein reines Einzelblock-Argument die DC schließen kann ---
das ist keine Beweislücke, sondern ein struktureller Satz über die Grenzen
des Ansatzes.

\medskip
Die lineare DC-Schranke $\nu \gtrsim 0{,}585\,k$ ist
\textbf{für die extremale Familie unerreichbar}, denn $\nu = O(\log k)$.
Dieser Befund ist exakt und widerlegungsfest.
\end{tcolorbox}

\subsection{Vergleichstabelle: Nettofaktor für $m = 3$ ($r = 1$)}

\begin{center}
\renewcommand{\arraystretch}{1.4}
\begin{tabular}{ccccc}
\toprule
$k$ & $\nu = \vii(3n_k+1)$ & Nettofaktor $\approx$ & Bedarf~$\nu$ & Kontrahiert? \\
\midrule
1 & $3 + \vii(3) = 4$ & $(3/2)^2/2^3 = 9/32 \approx 0{,}28$ & $>2{,}17$ & Ja \\
2 & $2$ (da $k+r+1=5$ ung.) & $(3/2)^3/2^1 = 27/8 \approx 3{,}4$ & $>2{,}75$ & Nein \\
3 & $3+\vii(4)=5$ & $(3/2)^4/2^3 = 81/128 \approx 0{,}63$ & $>3{,}34$ & Nein \\
4 & $2$ (da $k+r+1=7$ ung.) & $(3/2)^5/2^0 = 243/32 \approx 7{,}6$ & $>3{,}92$ & Nein \\
6 & $2$ (da $k+r+1=9$ ung.) & $(3/2)^7/2^0 \approx 17{,}1$ & $>5{,}1$ & Nein \\
7 & $3+\vii(8)=6$ & $(3/2)^8/2^4 = 6561/4096 \approx 1{,}6$ & $>5{,}68$ & Nein \\
10 & $2$ & $(3/2)^{11}/2^0 \approx 86$ & $>7{,}4$ & Nein \\
\bottomrule
\end{tabular}
\end{center}

Für $k \geq 3$ kontrahiert der einzelne $\mathrm{C}^k\mathrm{A}$-Block
(mit $m = 3$) selbst im besten verfügbaren Fall ($\nu = 5$) nicht mehr.
Der Orbit muss weitere Blöcke durchlaufen.

% ====================================================
\section{Die Endpunkt-Familie $2\cdot 3^N - 1$}
\label{sec:endpunkte}
% ====================================================

\subsection{Herkunft und Bedeutung}

Aus dem vorigen Abschnitt folgt: Die Endpunkte extremaler C-Ketten
(mit $n_0+1 = 2^{k+1}\cdot 3^r$) liegen alle auf der Familie
\[
  n_k = 2\cdot 3^{k+r} - 1 =: 2\cdot 3^N - 1, \quad N = k+r.
\]
Dies ist eine \emph{einparametrige Familie}, deren Verhalten unter der
Collatz-Abbildung vollständig durch LTE klassifizierbar ist.

\subsection{Vollständige LTE-Klassifikation}

\begin{proposition}[Klassifikation der Familie $2\cdot 3^N - 1$]
\label{prop:endpunktfamilie}
Für $n = 2\cdot 3^N - 1$ gilt:
\[
  \vii\bigl(3n + 1\bigr)
  = \vii\bigl(6\cdot 3^N - 2\bigr)
  = 1 + \vii(3^{N+1} - 1)
  =
  \begin{cases}
    2 & \text{falls $N+1$ ungerade (d.\,h.\ $N$ gerade),}\\[2pt]
    3 + \vii(N+1) & \text{falls $N+1$ gerade (d.\,h.\ $N$ ungerade).}
  \end{cases}
\]
Insbesondere: Für ungerades $N$ ist $\vii(3n+1)$ explizit und vollständig
bestimmt. Die Größe ist für diese Familie \textbf{nicht offen}.
\end{proposition}

\begin{beweis}
$3n + 1 = 3(2\cdot 3^N - 1) + 1 = 6\cdot 3^N - 2 = 2(3^{N+1} - 1)$.
Daher $\vii(3n+1) = 1 + \vii(3^{N+1}-1)$.
LTE für $p=2$: $\vii(3^M - 1) = 1$ falls $M$ ungerade, $\vii(M)+2$ falls
$M$ gerade.
\end{beweis}

\subsection{Explizite Wertetabelle}

\begin{center}
\renewcommand{\arraystretch}{1.4}
\begin{tabular}{ccccc}
\toprule
$N$ & $n = 2\cdot 3^N - 1$ & $N+1$ & $\vii(N+1)$ & $\vii(3n+1)$ \\
\midrule
1  & 5           & 2 (ger.)  & 1 & $3+1 = 4$ \\
2  & 17          & 3 (ung.)  & — & $2$ \\
3  & 53          & 4 (ger.)  & 2 & $3+2 = 5$ \\
4  & 161         & 5 (ung.)  & — & $2$ \\
5  & 485         & 6 (ger.)  & 1 & $3+1 = 4$ \\
6  & 1457        & 7 (ung.)  & — & $2$ \\
7  & 4373        & 8 (ger.)  & 3 & $3+3 = 6$ \\
8  & 13121       & 9 (ung.)  & — & $2$ \\
9  & 39365       & 10 (ger.) & 1 & $3+1 = 4$ \\
10 & 118097      & 11 (ung.) & — & $2$ \\
\bottomrule
\end{tabular}
\end{center}

\begin{tcolorbox}[ergebnis, title={Vollständige Klassifikation}]
Für $n = 2\cdot 3^N - 1$ ist $\vii(3n+1)$ für alle $N \in \N$ explizit
und geschlossen berechenbar. Das Muster:
\begin{itemize}
  \item $N$ gerade: $\vii(3n+1) = 2$ (minimale Tiefe, schlechtester Fall).
  \item $N$ ungerade: $\vii(3n+1) = 3 + \vii(N+1)$
    (Tiefe wächst logarithmisch mit $N$).
\end{itemize}
Diese Familie hat \emph{kontrollierbares, vollständig analysiertes Verhalten}.
\end{tcolorbox}

\subsection{Beobachtung: Kein Wachstum der Tiefe über $O(\log N)$}

Die Tabelle zeigt das strukturelle Muster: Die Tiefe $\vii(3n+1)$ ist für
$N$ ungerade genau $3 + \vii(N+1)$. Da $\vii(N+1) \leq \log_2(N+1)$, wächst
die Tiefe höchstens logarithmisch. Ein lineares Wachstum (wie für Kontraktion
nötig) tritt \emph{nie} auf. Die Endpunkte extremaler C-Ketten liefern also
systematisch zu flache A-Schritte für große~$k$.

% ====================================================
\section{Übergang zur Mehrblock-Dynamik}
\label{sec:mehrblock}
% ====================================================

\subsection{Definitionen}

\begin{definition}[Block-Zerlegung einer Collatz-Bahn]
Eine \emph{Block-Zerlegung} der Bahn $(n_0, \U(n_0), \U^2(n_0), \ldots)$
ist die Einteilung in maximale Teilfolgen vom Typ
\[
  B_j = \mathrm{C}^{l_j}\!\cdot\mathrm{A},
\]
wobei $l_j \geq 0$ die Länge der $j$-ten C-Kette und der abschließende
A-Schritt obligatorisch ist (nach Thm.~3.3 in \textit{collatz\_schlussartikel.tex}).
Der \emph{Nettofaktor} des $j$-ten Blocks ist:
\[
  F(B_j) := \frac{n_0^{(j+1)}}{n_0^{(j)}}
  \approx \frac{(3/2)^{l_j+1}}{2^{\vii(3n_{l_j}^{(j)}+1)}},
\]
wobei $n_0^{(j)}$ der Startwert des $j$-ten Blocks ist.
\end{definition}

\subsection{Die Mehrblock-Bedingung}

\begin{tcolorbox}[kernidee, title={Das richtige Beweisziel}]
Die Collatz-Vermutung ist äquivalent zur Aussage: Für jedes $n_0 \in \N$
gilt
\[
  \prod_{j=1}^{J} F(B_j) \;\xrightarrow{J \to \infty}\; 0.
\]
In logarithmischer Form: Das \emph{Mehrblock-Drift-Kriterium}
\[
  \limsup_{J\to\infty} \frac{1}{J} \sum_{j=1}^{J} \log F(B_j) < 0
\]
ist notwendig und hinreichend für die Beschränktheit der Bahn von $n_0$.
\end{tcolorbox}

\begin{bemerkung}[Vergleich mit Einzelblock-Kriterium]
Das Einzelblock-Kriterium verlangt $F(B_j) < 1$ für alle $j$. Das
Mehrblock-Kriterium verlangt nur, dass der \emph{Cesàro-Mittelwert}
der Logarithmen negativ ist. Letzteres ist wesentlich schwächer und
prinzipiell erreichbar --- auch wenn einzelne Blöcke $F(B_j) \gg 1$ haben.
\end{bemerkung}

\subsection{Strukturelles Bild}

Das Mehrblock-Drift-Kriterium besagt: Für jeden Block mit großem
Nettofaktor ($l_j$ groß, $\nu_j$ klein) müssen nachfolgende Blöcke
mit kleinem Nettofaktor ($\nu_j$ groß) ausgleichen. Dies erfordert
\emph{keine} Kontraktion in jedem Schritt, sondern nur eine
\emph{negative mittlere Steigung} im logarithmischen Bild.

\begin{tcolorbox}[lte, title={Mittlerer Nettofaktor nach LTE}]
Modellrechnung: Sei die C-Kettenlänge gleichverteilt in
$\{0, 1, \ldots, L\}$ und $\nu \in \{2, 3\}$ mit gleicher Wahrscheinlichkeit.
Dann ist der erwartete Logarithmus des Nettofaktors:
\[
  \mathbb{E}[\log F(B)] = \frac{1}{L+1}\sum_{l=0}^{L} \Bigl[
    (l+1)\log\tfrac{3}{2} - \bar\nu \log 2
  \Bigr]
  = \tfrac{L+2}{2} \log\tfrac{3}{2} - \bar\nu\log 2,
\]
wobei $\bar\nu \approx 2{,}5$. Für $L \leq L^*$ mit
$\frac{L^*+2}{2}\log\frac{3}{2} < \bar\nu\log 2$, d.\,h.\
$L^* < \frac{2\bar\nu\log 2}{\log(3/2)} - 2 \approx 6{,}4$, ist der
mittlere Nettofaktor $< 1$. Dies stützt die Konvergenz für \emph{typische}
Bahnen, beweist sie aber nicht für alle.
\end{tcolorbox}

% ====================================================
\section{Präzise Formulierung des Beweisziels}
\label{sec:beweisziel}
% ====================================================

\subsection{Die richtige Form der DC}

\begin{tcolorbox}[kernidee, title={Mehrblock-DC (richtige Form)}]
\textbf{Ziel:} Zeige für jedes $n_0 \in \N$:
\[
  \limsup_{J\to\infty} \frac{1}{J} \sum_{j=1}^{J} \log F(B_j) < 0.
\]
Dies ist äquivalent zur Collatz-Vermutung (für die Klasse der divergenzfreien
Bahnen), aber in einer Form, die \emph{ergodisch-theoretisch angreifbar} ist:
Es geht um das Vorzeichen eines Lyapunov-Exponenten.
\end{tcolorbox}

\subsection{Drei Präzisierungsstufen}

\begin{enumerate}[label=\textbf{(\arabic*)}, leftmargin=*]
  \item \textbf{Fast alle Bahnen (Tao 2022):}
    Tao zeigt: Für fast alle $n_0$ (im Sinne der logarithmischen Dichte)
    nehmen Orbits fast beschränkte Werte an. Das entspricht einer
    ``fast überall''-Version des Drift-Kriteriums.

  \item \textbf{Alle Bahnen mit beschränkter Kettenlänge:}
    Falls $l_j \leq C$ für alle $j$ (gleichmäßig beschränkte Kettenlängen),
    dann gilt $F(B_j) \leq (3/2)^{C+1}/2 < \infty$. Das Drift-Kriterium
    reduziert sich auf eine endliche Zahl von Fällen und ist analytisch
    handhabbar.

  \item \textbf{Alle Bahnen (Collatz-Vermutung):}
    Das allgemeine Drift-Kriterium für unbeschränkte $l_j$ erfordert globale
    Methoden. Die LTE-Analyse zeigt: Für große $l_j$ sind die Nettofaktoren
    groß ($\gg 1$), aber selten (Dichte $\leq 2^{-(l_j+1)}$). Die Frage
    ist, ob Seltenheit + logarithmisches Mittel negativ erzwingen.
\end{enumerate}

\subsection{Verbindung zur Ergodizität}

Die Drift-Bedingung lässt sich als Bedingung an den \emph{Lyapunov-Exponenten}
des Transferoperators $\mathcal{T}$ der Collatz-Abbildung formulieren:
\[
  \lambda := \lim_{J\to\infty} \frac{1}{J}\sum_{j=1}^J \log F(B_j)
  = \log\frac{3}{2} - \bar\nu\,\log 2 \stackrel{?}{<} 0.
\]
Die Frage $\bar\nu > \log_2 3 \approx 1{,}585$ (mittlere EA-Tiefe über
$\log_2 3$) ist äquivalent zur Collatz-Vermutung in dieser Formulierung.

\begin{tcolorbox}[konj, title={Drift-Konjektur}]
Für jede Collatz-Bahn gilt:
\[
  \liminf_{J\to\infty} \frac{1}{J} \sum_{j=1}^J \vii(3n_{l_j}^{(j)}+1) > \log_2 3.
\]
Das ist die \emph{rechte} Form der Durchlaufbedingung.
\end{tcolorbox}

% ====================================================
\section{Verbindung zur $E_8$-Analyse}
\label{sec:e8}
% ====================================================

\subsection{Return-Map-Argument}

In \textit{collatz\_oktonionen\_beweis.tex} wird ein $E_8$-Normabstieg
nachgewiesen, der bei jedem EA-Schritt für hinreichend große~$n$ eintritt.
Das Mehrblock-Drift-Argument ist die \emph{arithmetische Parallele} dazu:

\begin{center}
\renewcommand{\arraystretch}{1.4}
\begin{tabular}{p{4.5cm}p{1.5cm}p{5.5cm}}
\toprule
\textbf{$E_8$-Analyse} & & \textbf{Mehrblock-Drift} \\
\midrule
$E_8$-Norm fällt nicht nach jedem Schritt &$\longleftrightarrow$&
  Nettofaktor $F(B_j)$ ist nicht immer $<1$ \\[4pt]
Aber: Norm fällt nach jedem EA-Schritt &$\longleftrightarrow$&
  Einzelne große Blöcke werden durch folgende kompensiert \\[4pt]
Return-Map: Norm fällt nach jedem Rückzug in die ``Normalzone'' &$\longleftrightarrow$&
  Cesàro-Mittel von $\log F(B_j)$ ist negativ \\
\bottomrule
\end{tabular}
\end{center}

\subsection{Strukturelle Konsequenz}

\begin{tcolorbox}[kernidee, title={Gemeinsamer Kern}]
Sowohl das $E_8$-Argument als auch das Mehrblock-Drift-Argument zeigen:
\textbf{Lokale Kontraktion ist nicht das Ziel}. Beide Methoden liefern
Aussagen über \emph{akkumuliertes} Verhalten, nicht über einzelne Schritte.

Der $E_8$-Normabstieg bei EA-Schritten und das negative Cesàro-Mittel
von $\log F(B_j)$ sind zwei Seiten derselben strukturellen Aussage:
Die Collatz-Dynamik \emph{driftet} im Mittel nach unten.

Was fehlt, ist der Nachweis, dass dieser Mittelwert auch
\emph{für jeden einzelnen Orbit} negativ ist --- nicht nur im
statistischen Durchschnitt.
\end{tcolorbox}

\subsection{Formulierung als offenes Problem}

\begin{tcolorbox}[lücke, title={Präzise verbleibende Lücke}]
Die einzige verbleibende Lücke lautet jetzt:

\medskip
\begin{center}
\textit{Zeige: Für jedes $n_0 \in \N$ gilt\\[6pt]
$\displaystyle\limsup_{J\to\infty}\frac{1}{J}\sum_{j=1}^{J}\log F(B_j) < 0.$}
\end{center}

\medskip
Das ist der \emph{Lyapunov-Exponent der Collatz-Dynamik}.
Sein Vorzeichen ist das Collatz-Problem in seiner schärfsten Form.
\end{tcolorbox}

% ====================================================
% Literatur
% ====================================================
\begin{thebibliography}{9}

\bibitem{hoffbauer-schluss}
T.\ Hoffbauer,
\textit{EABC-Rahmen und Collatz-Dynamik: Ein strukturierter Beweisversuch},
Mathematische Texte, \texttt{collatz\_schlussartikel.tex} (2026).

\bibitem{hoffbauer-c-ketten}
T.\ Hoffbauer,
\textit{C-Ketten in der Collatz-Dynamik: 2-adische Analyse},
Mathematische Texte, \texttt{collatz\_c\_ketten\_2adisch.tex} (2025/2026).

\bibitem{hoffbauer-5glatt}
T.\ Hoffbauer,
\textit{C-Ketten, 5-glatte Zahlen und Abstiegsstruktur},
Mathematische Texte, \texttt{collatz\_5glatt\_dc.tex} (2026).

\bibitem{hoffbauer-oktonionen}
T.\ Hoffbauer,
\textit{Collatz-Dynamik auf ganzzahligen Oktonionen: Beweisversuch über das $E_8$-Gitter},
Mathematische Texte, \texttt{collatz\_oktonionen\_beweis.tex} (2025).

\bibitem{tao-collatz}
T.\ Tao,
\textit{Almost all orbits of the Collatz map attain almost bounded values},
Forum of Mathematics Pi \textbf{10} (2022), e12.

\bibitem{lte}
P.\ Kunec / V.\ Pop / B.\ Eskenazi et al.,
\textit{Lifting the Exponent Lemma (LTE)},
Art of Problem Solving, Community Notes (2009--2024).

\end{thebibliography}

\end{document}
